%=============================================================
%  ESENCIA SEVILLA — DOCUMENTACIÓN TÉCNICA COMPLETA
%  Versión 1.0.0  |  Abril 2026
%=============================================================
\documentclass[12pt,a4paper,twoside,openright]{report}

%--- Codificación y lengua --------------------------------
\usepackage[utf8]{inputenc}
\usepackage[T1]{fontenc}
\usepackage[spanish,es-nolayout]{babel}
\usepackage{lmodern}

%--- Geometría de página ----------------------------------
\usepackage[
  top=3cm, bottom=3cm,
  inner=3.5cm, outer=2.5cm,
  headheight=15pt,
  footskip=1.5cm
]{geometry}

%--- Tipografía -------------------------------------------
\usepackage{microtype}
\usepackage{setspace}
\onehalfspacing
\usepackage{parskip}

%--- Paleta de colores del diseño ------------------------
\usepackage[dvipsnames,svgnames,table]{xcolor}
\definecolor{Terracota}{HTML}{C25A3A}
\definecolor{TerracotaClaro}{HTML}{FAE5DC}
\definecolor{TerracotaOscuro}{HTML}{82321B}
\definecolor{Azulejo}{HTML}{2A5A8C}
\definecolor{AzulejoClaro}{HTML}{EFF5FB}
\definecolor{Ocre}{HTML}{D9A760}
\definecolor{Crema}{HTML}{F7F0E3}
\definecolor{Tinta}{HTML}{2B1E15}
\definecolor{GrisClaro}{HTML}{F4F4F4}
\definecolor{GrisMedio}{HTML}{CCCCCC}
\definecolor{Verde}{HTML}{16A34A}
\definecolor{Naranja}{HTML}{D97706}

%--- Encabezados y pies de página ------------------------
\usepackage{fancyhdr}
\pagestyle{fancy}
\fancyhf{}
\fancyhead[LE]{\small\textcolor{Terracota}{\textbf{Esencia Sevilla}} \textcolor{GrisMedio}{|} Documentación Técnica v1.0}
\fancyhead[RO]{\small\textcolor{Tinta}{\leftmark}}
\fancyfoot[LE,RO]{\small\textcolor{GrisMedio}{\thepage}}
\renewcommand{\headrulewidth}{0.4pt}
\renewcommand{\headrule}{\hbox to\headwidth{\color{TerracotaClaro}\leaders\hrule height \headrulewidth\hfill}}

%--- Formato de capítulos y secciones --------------------
\usepackage{titlesec}
\titleformat{\chapter}[display]
  {\normalfont\Huge\bfseries\color{Tinta}}
  {{\Large\color{Terracota}\chaptertitlename\ \thechapter}}
  {16pt}
  {\Huge}[\vspace{4pt}{\color{TerracotaClaro}\hrule height 2pt}]

\titleformat{\section}
  {\normalfont\Large\bfseries\color{Azulejo}}
  {\thesection}{1em}{}[\vspace{2pt}{\color{AzulejoClaro}\hrule height 0.5pt}]

\titleformat{\subsection}
  {\normalfont\large\bfseries\color{Terracota}}
  {\thesubsection}{1em}{}

\titleformat{\subsubsection}
  {\normalfont\normalsize\bfseries\color{Tinta}}
  {\thesubsubsection}{1em}{}

%--- Tablas -----------------------------------------------
\usepackage{booktabs}
\usepackage{longtable}
\usepackage{tabularx}
\usepackage{multirow}
\usepackage{array}
\usepackage{colortbl}
\newcolumntype{L}[1]{>{\raggedright\arraybackslash}p{#1}}
\newcolumntype{C}[1]{>{\centering\arraybackslash}p{#1}}
\newcolumntype{R}[1]{>{\raggedleft\arraybackslash}p{#1}}
\newcolumntype{X}{>{\raggedright\arraybackslash}X}

%--- Código fuente ----------------------------------------
\usepackage{listings}
\usepackage{lstautogobble}

\lstdefinelanguage{TypeScript}{
  keywords={const,let,var,function,return,if,else,while,for,
            export,import,from,async,await,typeof,interface,
            type,class,extends,implements,new,this,null,
            undefined,true,false,void,string,number,boolean},
  morecomment=[l]{//},
  morecomment=[s]{/*}{*/},
  morestring=[b]",
  morestring=[b]',
  morestring=[b]`,
  sensitive=true
}

\lstdefinestyle{tsStyle}{
  language=TypeScript,
  backgroundcolor=\color{GrisClaro},
  basicstyle=\ttfamily\footnotesize,
  keywordstyle=\color{Azulejo}\bfseries,
  stringstyle=\color{Terracota},
  commentstyle=\color{gray!70}\itshape,
  numberstyle=\tiny\color{GrisMedio},
  numbers=left,
  numbersep=8pt,
  breaklines=true,
  breakatwhitespace=true,
  columns=flexible,
  keepspaces=true,
  showstringspaces=false,
  frame=leftline,
  framesep=8pt,
  rulecolor=\color{Terracota},
  xleftmargin=20pt,
  xrightmargin=4pt,
  autogobble=true,
  tabsize=2,
}

\lstdefinestyle{sqlStyle}{
  language=SQL,
  backgroundcolor=\color{AzulejoClaro},
  basicstyle=\ttfamily\footnotesize,
  keywordstyle=\color{Azulejo}\bfseries,
  stringstyle=\color{Terracota},
  commentstyle=\color{gray!70}\itshape,
  numberstyle=\tiny\color{GrisMedio},
  numbers=left,
  numbersep=8pt,
  breaklines=true,
  frame=leftline,
  framesep=8pt,
  rulecolor=\color{Azulejo},
  xleftmargin=20pt,
  autogobble=true,
}

\lstdefinestyle{bashStyle}{
  language=bash,
  backgroundcolor=\color{Tinta!8},
  basicstyle=\ttfamily\footnotesize\color{Tinta},
  keywordstyle=\color{Terracota}\bfseries,
  stringstyle=\color{Verde!80!black},
  commentstyle=\color{gray!70}\itshape,
  numbers=none,
  breaklines=true,
  frame=none,
  xleftmargin=12pt,
  autogobble=true,
}

\lstdefinestyle{jsonStyle}{
  language=bash,
  backgroundcolor=\color{GrisClaro},
  basicstyle=\ttfamily\footnotesize,
  numbers=none,
  breaklines=true,
  frame=single,
  framesep=6pt,
  rulecolor=\color{GrisMedio},
  xleftmargin=8pt,
  xrightmargin=8pt,
  autogobble=true,
}

\lstset{style=tsStyle}

%--- Listas -----------------------------------------------
\usepackage{enumitem}
\setlist[itemize]{leftmargin=1.5em, itemsep=3pt, topsep=4pt}
\setlist[enumerate]{leftmargin=1.5em, itemsep=3pt, topsep=4pt}

%--- Cajas decorativas ------------------------------------
\usepackage{tcolorbox}
\tcbuselibrary{skins, breakable, listings}

\newtcolorbox{infobox}[2][]{
  enhanced, breakable,
  colback=TerracotaClaro!40,
  colframe=Terracota,
  coltitle=white,
  fonttitle=\bfseries\small,
  title=#2,
  left=8pt, right=8pt, top=6pt, bottom=6pt,
  arc=5pt,
  drop shadow={Terracota!30},
  #1
}

\newtcolorbox{warnbox}[2][]{
  enhanced, breakable,
  colback=Ocre!15,
  colframe=Ocre,
  coltitle=Tinta,
  fonttitle=\bfseries\small,
  title={⚠ #2},
  left=8pt, right=8pt, top=6pt, bottom=6pt,
  arc=5pt,
  #1
}

\newtcolorbox{successbox}[2][]{
  enhanced, breakable,
  colback=Verde!8,
  colframe=Verde,
  coltitle=white,
  fonttitle=\bfseries\small,
  title={✓ #2},
  left=8pt, right=8pt, top=6pt, bottom=6pt,
  arc=5pt,
  #1
}

\newtcolorbox{codebox}[2][]{
  enhanced, breakable,
  colback=GrisClaro,
  colframe=GrisMedio,
  coltitle=Tinta,
  fonttitle=\ttfamily\bfseries\small,
  title=#2,
  left=6pt, right=6pt, top=6pt, bottom=6pt,
  arc=3pt,
  #1
}

%--- Gráficos y diagramas ---------------------------------
\usepackage{graphicx}
\usepackage{float}
\usepackage{tikz}
\usetikzlibrary{shapes.geometric, arrows.meta, positioning, backgrounds,
                fit, shadows.blur, calc}

%--- Hiperlinks -------------------------------------------
\usepackage[
  colorlinks=true,
  linkcolor=Azulejo,
  urlcolor=Terracota,
  citecolor=Tinta,
  filecolor=Verde
]{hyperref}
\usepackage{url}

%--- Otros ------------------------------------------------
\usepackage{eurosym}
\usepackage{amssymb}
\usepackage{pifont}
\newcommand{\cmark}{\textcolor{Verde}{\ding{51}}}
\newcommand{\xmark}{\textcolor{red}{\ding{55}}}
\newcommand{\pending}{\textcolor{Naranja}{\textbf{!}}}

%--- Metadatos PDF ----------------------------------------
\hypersetup{
  pdftitle={Esencia Sevilla — Documentación Técnica v1.0.0},
  pdfauthor={Esencia Sevilla Dev Team},
  pdfsubject={Documentación técnica completa de la aplicación web},
  pdfkeywords={Next.js 14, Supabase, Stripe, TypeScript, i18n, PWA, Sevilla},
  pdfcreator={LaTeX + hyperref},
  pdfproducer={pdfTeX}
}

%=============================================================
\begin{document}
%=============================================================

\frontmatter

%--- Portada ------------------------------------------------
\begin{titlepage}
\pagecolor{Crema}
\afterpage{\nopagecolor}
\centering
\vspace*{1.5cm}

{\color{Terracota}\rule{\textwidth}{4pt}}
\vspace{1.5cm}

{\fontsize{42}{52}\selectfont\bfseries\color{Tinta}Esencia Sevilla}

\vspace{0.6cm}
{\Large\color{Terracota}\itshape Imaginero Luis Alvarez Duarte N7 · 41008 Sevilla}

\vspace{0.3cm}
{\normalsize\color{Tinta!60} Nº Registro Turístico AT/SE/03584}

\vspace{1.5cm}
{\color{GrisMedio}\rule{\textwidth}{0.5pt}}
\vspace{1cm}

{\fontsize{28}{36}\selectfont\bfseries\color{Azulejo}Documentación Técnica}

\vspace{0.5cm}
{\large\color{Tinta!80} Aplicación Web de Gestión de Apartamento Turístico}

\vspace{0.4cm}
{\normalsize\color{Tinta!60}
Next.js 14 · TypeScript · Tailwind CSS · Supabase · Stripe · OpenRouter
}

\vspace{1.5cm}
{\color{GrisMedio}\rule{0.6\textwidth}{0.5pt}}
\vspace{1cm}

\begin{tabular}{@{}ll@{}}
  \textcolor{Tinta!50}{\small Versión} & \textbf{1.0.0} \\[6pt]
  \textcolor{Tinta!50}{\small Fecha} & \textbf{Abril 2026} \\[6pt]
  \textcolor{Tinta!50}{\small Repositorio} & \texttt{\small github.com/danimap27/esencia-sevilla} \\[6pt]
  \textcolor{Tinta!50}{\small Licencia} & \textbf{Privado} \\
\end{tabular}

\vfill
{\color{Terracota}\rule{\textwidth}{4pt}}
\end{titlepage}

%--- Página en blanco y tabla de contenidos ---------------
\cleardoublepage
\tableofcontents
\clearpage
\listoftables
\clearpage
\listoffigures

\mainmatter

%=============================================================
\chapter{Descripción General del Producto}
%=============================================================

\section{Visión y objetivos de negocio}

\textbf{Esencia Sevilla} es una aplicación web de producción para la comercialización
directa y la gestión integral de un apartamento turístico de alta gama ubicado en el
corazón histórico de Sevilla, a menos de 10 minutos a pie de la Catedral, el Real Alcázar
y la Torre del Oro.

La plataforma nace para resolver el principal problema de los propietarios de alquiler
vacacional: la \textbf{dependencia total de las OTAs} (Online Travel Agencies) como
Booking.com y Airbnb, que cobran comisiones de entre el 15\% y el 25\% en cada reserva.
La estrategia adoptada es ofrecer al viajero un \textbf{descuento permanente del 10\%}
sobre el precio publicado en las OTAs, de manera que el huésped pague menos \emph{y}
el propietario ingrese más, eliminando el intermediario en ambas direcciones.

\begin{infobox}{Datos del inmueble}
\begin{tabular}{@{}L{4cm}L{7.5cm}@{}}
  \textbf{Nombre:}         & Esencia Sevilla \\[3pt]
  \textbf{Dirección:}      & Imaginero Luis Alvarez Duarte N7, 41008 Sevilla \\[3pt]
  \textbf{Coordenadas:}    & 37.38862°N,\ \ −5.98230°W \\[3pt]
  \textbf{Reg. Turístico:} & AT/SE/03584 (Junta de Andalucía) \\[3pt]
  \textbf{Capacidad:}      & 4 huéspedes · 2 dormitorios · 3 camas · 65 m² \\[3pt]
  \textbf{Precio base:}    & \euro{}142/noche (temporada media) \\[3pt]
  \textbf{Tarifa limpieza:}& \euro{}60 única por estancia \\[3pt]
  \textbf{Tasa turística:} & \euro{}1,50 por persona/noche \\[3pt]
  \textbf{Estancia mínima:}& 2 noches \\
\end{tabular}
\end{infobox}

\subsection{Objetivos de negocio}

\begin{longtable}{L{4.5cm} L{8cm}}
  \toprule
  \rowcolor{AzulejoClaro}
  \textbf{Objetivo} & \textbf{Implementación técnica} \\
  \midrule
  Captación reservas directas & Sistema propio con Stripe Checkout, −10\% vs OTA \\
  \addlinespace
  Reducir dependencia de OTAs & Sincronización iCal bidireccional + bloqueo automático de fechas \\
  \addlinespace
  Presencia internacional & 6 idiomas: ES, EN, FR, DE, IT, PT (next-intl) \\
  \addlinespace
  Posicionamiento SEO orgánico & JSON-LD, hreflang, sitemap dinámico, blog de Sevilla \\
  \addlinespace
  Experiencia superior al huésped & Chatbot IA 24/7, guía turística interactiva, pre-check-in digital \\
  \addlinespace
  Cumplimiento legal & RGPD, SES.HOSPEDERÍA, normativa turística andaluza \\
  \bottomrule
\end{longtable}

\section{Puntos de interés cercanos}

\begin{longtable}{L{5cm} C{3cm}}
  \toprule
  \textbf{Monumento / Lugar} & \textbf{Distancia a pie} \\
  \midrule
  Barrio de Santa Cruz & 6 min \\
  Torre del Oro & 7 min \\
  Catedral de Sevilla & 8 min \\
  La Giralda & 9 min \\
  Archivo de Indias & 9 min \\
  Real Alcázar & 10 min \\
  Plaza de España & 12 min \\
  Museo de Bellas Artes & 15 min \\
  \bottomrule
\end{longtable}

%=============================================================
\chapter{Stack Tecnológico}
%=============================================================

\section{Frontend}

\begin{longtable}{L{3.5cm} C{2cm} L{7cm}}
  \toprule
  \rowcolor{AzulejoClaro}
  \textbf{Librería} & \textbf{Versión} & \textbf{Descripción} \\
  \midrule
  Next.js & 14.2.5 & Framework principal con App Router, RSC, Route Handlers \\
  React & 18.3.1 & UI library con Suspense y Concurrent Rendering \\
  TypeScript & \textasciicircum 5 & Tipado estático en toda la aplicación \\
  Tailwind CSS & \textasciicircum 3.4.6 & Utilidades CSS con design tokens personalizados \\
  next-intl & \textasciicircum 3.17.3 & i18n nativo para App Router con Server Components \\
  framer-motion & \textasciicircum 11.2.12 & Animaciones declarativas de alta calidad \\
  lucide-react & \textasciicircum 0.411.0 & Librería de 1.000+ iconos SVG consistentes \\
  react-day-picker & \textasciicircum 8.10.1 & Selector de fechas con modo DateRange y localización \\
  react-leaflet & \textasciicircum 4.2.1 & Wrapper React para Leaflet (carga con SSR desactivado) \\
  leaflet & \textasciicircum 1.9.4 & Motor de mapas OpenStreetMap \\
  react-hot-toast & \textasciicircum 2.4.1 & Notificaciones toast no intrusivas \\
  Radix UI & \textasciicircum 1.x & Primitivas UI accesibles: Dialog, Accordion, Tabs, Select \\
  date-fns & \textasciicircum 3.6.0 & Manipulación de fechas con soporte de 6 locales \\
  clsx + tailwind-merge & --- & Composición segura de clases CSS \\
  \bottomrule
\end{longtable}

\section{Backend y servicios externos}

\begin{longtable}{L{3.5cm} C{2cm} L{7cm}}
  \toprule
  \rowcolor{AzulejoClaro}
  \textbf{Servicio} & \textbf{Versión} & \textbf{Descripción} \\
  \midrule
  Supabase & \textasciicircum 2.44.4 & PostgreSQL gestionado + RLS + API REST automática \\
  Stripe & \textasciicircum 16.2.0 & Checkout Sessions, webhooks, PCI-DSS Level 1 \\
  @stripe/stripe-js & \textasciicircum 4.1.0 & Stripe.js: Apple Pay, Google Pay, Payment Request \\
  Resend & \textasciicircum 3.4.0 & Emails transaccionales HTML con dominio propio \\
  OpenRouter & REST & Proxy de IA: Claude, GPT-4, Mistral... \\
  Twilio & REST & WhatsApp Business API (notificaciones propietario) \\
  ical.js & \textasciicircum 2.0.1 & Parser estándar iCalendar RFC 5545 \\
  \bottomrule
\end{longtable}

\section{Herramientas de desarrollo y build}

\begin{longtable}{L{4cm} C{2.5cm} L{6cm}}
  \toprule
  \rowcolor{AzulejoClaro}
  \textbf{Herramienta} & \textbf{Versión} & \textbf{Uso} \\
  \midrule
  Zod & \textasciicircum 3.23.8 & Validación de esquemas en todas las API routes \\
  qrcode & \textasciicircum 1.5.4 & Generación de QR en el panel de administración \\
  next-pwa & \textasciicircum 5.6.0 & Service Worker + Workbox para PWA \\
  ESLint & \textasciicircum 8 & Linting con reglas de Next.js \\
  PostCSS & \textasciicircum 8.4.39 & Autoprefixer + pipeline Tailwind \\
  \bottomrule
\end{longtable}

%=============================================================
\chapter{Arquitectura del Sistema}
%=============================================================

\section{Diagrama general}

\begin{figure}[H]
\centering
\resizebox{\textwidth}{!}{%
\begin{tikzpicture}[
  node distance=1.5cm and 2cm,
  >=Stealth,
  client/.style={rectangle, rounded corners=6pt, draw=Azulejo, fill=AzulejoClaro,
                 minimum width=12cm, minimum height=1.2cm, font=\small\bfseries, text=Azulejo},
  server/.style={rectangle, rounded corners=4pt, draw=Terracota, fill=TerracotaClaro!50,
                 minimum width=5cm, minimum height=1.1cm, font=\small, text centered},
  ext/.style={rectangle, rounded corners=4pt, draw=Tinta!40, fill=Crema,
              minimum width=3cm, minimum height=1cm, font=\small, text centered},
  arr/.style={->, thick, color=Tinta!50},
  label/.style={font=\scriptsize\itshape, text=Tinta!60},
]
  \node[client] (browser) {Browser / PWA — React 18 + next-intl + Leaflet + DayPicker};

  \node[server, below=1.5cm of browser, xshift=-3.2cm] (rsc) {Server Components\\(pages, sitemap, RSC)};
  \node[server, below=1.5cm of browser, xshift=3.2cm]  (api)  {API Routes\\(/api/*)};

  \node[ext, below=2cm of rsc, xshift=-2.5cm]  (supabase)    {Supabase\\PostgreSQL};
  \node[ext, below=2cm of api, xshift=-1cm]    (stripe)      {Stripe\\Checkout};
  \node[ext, below=2cm of api, xshift=2.5cm]   (openrouter)  {OpenRouter\\(IA)};
  \node[ext, below=3.5cm of rsc, xshift=2.5cm] (resend)      {Resend\\+ Twilio};
  \node[ext, below=5cm of supabase]             (ical)        {Booking.com\\Airbnb (iCal)};

  % Caja Next.js
  \begin{pgfonlayer}{background}
    \node[draw=Terracota, dashed, rounded corners=8pt, fill=TerracotaClaro!20,
          fit=(rsc)(api), inner sep=15pt,
          label={[color=Terracota,font=\small\bfseries]above:Next.js 14 App Router}] {};
  \end{pgfonlayer}

  % Flechas
  \draw[arr] (browser) -- (rsc);
  \draw[arr] (browser) -- (api);
  \draw[arr] (api) -- (supabase);
  \draw[arr] (api) -- (stripe);
  \draw[arr] (api) -- (openrouter);
  \draw[arr] (api) -- (resend);
  \draw[arr] (supabase) -- node[label, right]{iCal sync} (ical);

  % Etiqueta HTTPS
  \node[label, right=0.1cm of browser, yshift=-0.8cm] {HTTPS/2 · TLS 1.3};
\end{tikzpicture}%
}
\caption{Arquitectura general del sistema Esencia Sevilla}
\label{fig:architecture}
\end{figure}

\section{Patrones de renderizado}

\begin{longtable}{L{4cm} L{8.5cm}}
  \toprule
  \rowcolor{AzulejoClaro}
  \textbf{Patrón} & \textbf{Uso en la aplicación} \\
  \midrule
  Server Components (RSC) & Layout, Footer, páginas legales, sitemap, robots \\
  Client Components & Hero, BookingSection, Gallery, FAQ, Reviews, Admin \\
  Dynamic Import (SSR off) & MapSection (Leaflet), ChatWidget (window dependency) \\
  API Routes & Todas las interacciones con BD, Stripe, IA y auth \\
  Edge Middleware & Detección de locale y redirección i18n \\
  \bottomrule
\end{longtable}

\section{Flujo de una reserva}

\begin{enumerate}
  \item El usuario llega a la web; el \textbf{Middleware} detecta el locale (Accept-Language) y redirige a \texttt{/es} (o el locale detectado).
  \item \texttt{BookingSection} hace \texttt{GET /api/availability}: Supabase devuelve fechas bloqueadas + iCal de Booking.com y Airbnb. Caché de 5 min.
  \item El usuario selecciona fechas, huéspedes, extras y aplica cupón.
  \item Al pulsar ``Reservar directo'': \texttt{POST /api/bookings/checkout}.
  \begin{itemize}
    \item Zod valida el body de la petición.
    \item Se comprueba disponibilidad en Supabase (segunda validación server-side).
    \item Se calcula el precio en el servidor (inmune a tampering).
    \item Se crea una Stripe Checkout Session.
    \item Se inserta la reserva con \texttt{status: 'pending'}.
    \item Se devuelve \texttt{\{ url, bookingRef \}}.
  \end{itemize}
  \item El usuario es redirigido a la hosted page de Stripe para completar el pago.
  \item Stripe envía \texttt{checkout.session.completed} a \texttt{POST /api/webhooks/stripe}.
  \begin{itemize}
    \item Se verifica la firma HMAC-SHA256.
    \item Se actualiza la reserva: \texttt{status: 'confirmed'}.
    \item Se insertan todas las fechas de la estancia en \texttt{blocked\_dates}.
    \item Se envía email de confirmación HTML (Resend).
    \item Se notifica al propietario por WhatsApp (Twilio).
  \end{itemize}
  \item El usuario es redirigido a \texttt{/\{locale\}/reserva-confirmada?ref=...}
\end{enumerate}

%=============================================================
\chapter{Sistema de Diseño}
%=============================================================

\section{Paleta de colores}

Los colores evocan la arquitectura y cultura sevillana: el \textbf{terracota} de las
tejas mudéjares, el \textbf{azulejo} andaluz, el \textbf{ocre} de las fachadas y la
\textbf{crema} del encalado.

\begin{longtable}{C{2cm} L{2.2cm} C{3cm} L{5cm}}
  \toprule
  \textbf{Muestra} & \textbf{Token} & \textbf{Hex} & \textbf{Uso principal} \\
  \midrule
  \cellcolor{Terracota} & terracota-500 & \#C25A3A & CTAs, acentos, hover states \\
  \cellcolor{TerracotaClaro} & terracota-100 & \#FAE5DC & Fondos suaves, info boxes \\
  \cellcolor{Azulejo} & azulejo-500 & \#2A5A8C & Secciones, enlaces, información \\
  \cellcolor{AzulejoClaro} & azulejo-50 & \#EFF5FB & Fondos informativos \\
  \cellcolor{Ocre} & ocre-500 & \#D9A760 & Badges temporada alta, alertas \\
  \cellcolor{Crema} & crema & \#F7F0E3 & Fondo principal de la web \\
  \cellcolor{Tinta} & tinta & \#2B1E15 & Texto principal \\
  \bottomrule
\end{longtable}

\section{Tipografía}

\begin{longtable}{L{3cm} L{4cm} L{5.5cm}}
  \toprule
  \textbf{Fuente} & \textbf{Uso} & \textbf{Características} \\
  \midrule
  Cormorant Garamond & H1, H2, taglines & Serif elegante, cargada con next/font \\
  Inter & Cuerpo, UI, botones & Sans-serif legible, optical sizing \\
  \bottomrule
\end{longtable}

\section{Clases utilitarias de componente}

Definidas en \texttt{globals.css} para uso consistente en toda la aplicación:

\begin{longtable}{L{3.5cm} L{9cm}}
  \toprule
  \rowcolor{AzulejoClaro}
  \textbf{Clase CSS} & \textbf{Descripción} \\
  \midrule
  \texttt{.btn-primary} & Botón terracota sólido con hover, focus ring y transición \\
  \texttt{.btn-secondary} & Botón azulejo sólido \\
  \texttt{.btn-outline} & Botón con borde terracota, fondo transparente \\
  \texttt{.card} & Contenedor blanco con sombra suave y bordes redondeados \\
  \texttt{.section} & Padding vertical estándar (\texttt{py-16 sm:py-24}) \\
  \texttt{.input} & Input HTML estilizado con focus ring terracota \\
  \texttt{.badge} & Etiqueta pequeña con colores semánticos \\
  \texttt{.glass} & Fondo glass morfismo claro (backdrop-blur) \\
  \texttt{.glass-dark} & Fondo glass morfismo oscuro \\
  \texttt{.scrollbar-hidden} & Scroll sin barra visible \\
  \bottomrule
\end{longtable}

%=============================================================
\chapter{Componentes de Interfaz}
%=============================================================

\section{Header}

Barra de navegación sticky con efecto glass progresivo al superar 50 px de scroll
(umbral detectado con \texttt{useEffect} + \texttt{addEventListener('scroll')}).

En \textbf{desktop} presenta los links de sección con scroll suave y el selector de
idioma — un dropdown con las 6 banderas y nombres nativos. En \textbf{móvil},
un icono hamburguesa abre un overlay de pantalla completa con animación
\texttt{slideUpFade}.

El cambio de idioma usa \texttt{useRouter().replace()} con el nuevo locale,
preservando la ruta actual.

\section{Hero}

Sección de pantalla completa con cinco elementos:

\begin{enumerate}
  \item \textbf{Carousel automático} — 5 fotos con rotación cada 4.500 ms, indicadores
        de posición, gradiente negro para legibilidad.
  \item \textbf{Contador de ahorro en tiempo real} — calcula \texttt{calculatePrice()}
        para 3 noches/2 personas y muestra el ahorro vs Booking.com.
  \item \textbf{Badges de confianza} — valoración 4,9/5, +47 reseñas, WiFi ultra, auto check-in,
        nº de registro AT/SE/03584.
  \item \textbf{CTAs} — ``Reservar directo'' (scroll a \texttt{\#reservar}) y
        ``Ver en Booking.com'' (enlace externo de respaldo).
  \item \textbf{Landmarks cercanos} — 8 monumentos con emoji, nombre y tiempo a pie.
\end{enumerate}

\section{BookingSection}

El componente más complejo de la aplicación. Gestiona el flujo completo de reserva:

\begin{enumerate}
  \item \textbf{Tipo de reserva} — Pestaña Directa (−10\%, no reembolsable) / Booking.com.
  \item \textbf{Fechas} — \texttt{react-day-picker} en modo \texttt{DateRange}; las fechas
        bloqueadas se cargan de \texttt{GET /api/availability} al montar.
  \item \textbf{Huéspedes} — Contador 1–4. Recalcula la tasa turística en tiempo real.
  \item \textbf{Extras (Upsells)} — 5 opciones (traslado, welcome pack, late check-out,
        cuna de viaje, pack romántico).
  \item \textbf{Cupón} — Input con validación debounced de 800 ms contra
        \texttt{POST /api/bookings/validate-coupon}.
  \item \textbf{Desglose y pago} — Muestra el ahorro vs Booking.com y redirige a Stripe.
\end{enumerate}

\section{MapSection}

Cargado con \texttt{dynamic(\dots, \{ ssr: false \})} porque Leaflet requiere
el objeto \texttt{window}. Tiene 3 pestañas:

\begin{description}
  \item[Mapa] Marcador del apartamento (icono casa terracota) + 10 atracciones
              con colores por categoría. Clic en marcador → popup con descripción.
  \item[Rutas] 5 rutas turísticas con modal de detalle paso a paso.
  \item[Transporte] Tabla con líneas EA, C3, C4, 24 y Metro L1.
\end{description}

\section{ChatWidget}

Botón flotante (FAB) con indicador verde de disponibilidad. Al abrirse muestra la
ventana de chat con:
\begin{itemize}
  \item 4 chips de sugerencias predefinidas (localizados)
  \item Input de texto con envío por Enter o botón
  \item Indicador de escritura animado (3 puntos)
  \item Auto-scroll al último mensaje
\end{itemize}

Los mensajes se envían a \texttt{POST /api/chat}, que los reenvía a OpenRouter con
un system prompt de 50+ líneas con contexto completo del apartamento.

%=============================================================
\chapter{API REST — Referencia Completa}
%=============================================================

\section{GET /api/availability}

Devuelve las fechas bloqueadas para los próximos 120 días.

\begin{codebox}{Respuesta 200 — application/json}
\begin{lstlisting}[style=jsonStyle, numbers=none]
{
  "blockedDates": ["2026-05-10", "2026-05-11", "2026-06-01"]
}
\end{lstlisting}
\end{codebox}

\textbf{Lógica:} (1) Comprueba caché en memoria (TTL 5 min). (2) Consulta Supabase.
(3) Sincroniza iCal de Booking.com y Airbnb. (4) Deduplica y devuelve.

\section{POST /api/bookings/checkout}

Crea sesión Stripe y persiste la reserva como \texttt{pending}.

\begin{codebox}{Body requerido}
\begin{lstlisting}[style=jsonStyle, numbers=none]
{
  "checkIn":    "2026-06-15",
  "checkOut":   "2026-06-20",
  "guests":     2,
  "upsells":    [{ "id": "welcome-pack", "price": 25, "name": "Pack" }],
  "couponCode": "VERANO10",
  "locale":     "es"
}
\end{lstlisting}
\end{codebox}

\begin{longtable}{C{1.5cm} L{10cm}}
  \toprule
  \textbf{Código} & \textbf{Descripción} \\
  \midrule
  200 & \texttt{\{ "url": "https://checkout.stripe.com/...", "bookingRef": "ES-..." \}} \\
  400 & Validación Zod fallida o estancia < 2 noches \\
  409 & Fechas no disponibles (double-booking check) \\
  500 & Error interno \\
  \bottomrule
\end{longtable}

\section{POST /api/bookings/validate-coupon}

Valida un código sin consumir el uso. Verifica: \texttt{active}, \texttt{valid\_from},
\texttt{valid\_until}, \texttt{uses\_count < max\_uses}.

\begin{codebox}{Respuesta — cupón válido}
\begin{lstlisting}[style=jsonStyle, numbers=none]
{ "valid": true, "discountType": "percent", "discountValue": 10 }
\end{lstlisting}
\end{codebox}

\section{POST /api/chat}

Proxy a OpenRouter. Límite: 500 tokens de respuesta.

\section{POST /api/webhooks/stripe}

\begin{warnbox}{Solo para Stripe}
Esta ruta no puede llamarse manualmente. Requiere firma HMAC-SHA256 del payload.
\end{warnbox}

\begin{longtable}{L{6cm} L{7cm}}
  \toprule
  \rowcolor{AzulejoClaro}
  \textbf{Evento Stripe} & \textbf{Acciones en el servidor} \\
  \midrule
  \texttt{checkout.session.completed} &
    UPDATE booking → confirmed · INSERT blocked\_dates · Send email (Resend) · WhatsApp (Twilio) \\
  \addlinespace
  \texttt{checkout.session.expired} &
    UPDATE booking → cancelled \\
  \bottomrule
\end{longtable}

\section{Endpoints de administración}

\begin{longtable}{L{5.5cm} C{1.5cm} L{5.5cm}}
  \toprule
  \rowcolor{AzulejoClaro}
  \textbf{Ruta} & \textbf{Método} & \textbf{Acción} \\
  \midrule
  \texttt{/api/admin/auth} & POST & Login; cookie httpOnly 8h \\
  \texttt{/api/admin/bookings} & GET & Lista de reservas \\
  \texttt{/api/admin/block-date} & POST & Bloquear fecha manualmente \\
  \texttt{/api/admin/block-date} & DELETE & Desbloquear fecha manual \\
  \bottomrule
\end{longtable}

%=============================================================
\chapter{Base de Datos}
%=============================================================

\section{Modelo de datos}

El sistema usa \textbf{PostgreSQL} gestionado en Supabase con
\textbf{Row Level Security (RLS)} activo en las 4 tablas.

\subsection{Tabla bookings}

\begin{longtable}{L{4.2cm} L{3.5cm} L{5cm}}
  \toprule
  \rowcolor{AzulejoClaro}
  \textbf{Columna} & \textbf{Tipo} & \textbf{Descripción} \\
  \midrule
  \texttt{id} & UUID PK & \texttt{gen\_random\_uuid()} \\
  \texttt{booking\_ref} & TEXT UNIQUE & Ref. legible ES-XXXX-XXXX \\
  \texttt{check\_in} & DATE NOT NULL & Fecha de entrada \\
  \texttt{check\_out} & DATE NOT NULL & Fecha de salida \\
  \texttt{guests} & INTEGER & Nº de huéspedes (1--4) \\
  \texttt{guest\_name} & TEXT NOT NULL & Nombre completo \\
  \texttt{guest\_email} & TEXT NOT NULL & Email de contacto \\
  \texttt{total\_amount} & DECIMAL(10,2) & Total cobrado en \euro{} \\
  \texttt{status} & TEXT & pending/confirmed/cancelled/completed \\
  \texttt{payment\_intent\_id} & TEXT & ID de Stripe PaymentIntent \\
  \texttt{upsells} & JSONB & Array \texttt{[\{id, name, price\}]} \\
  \texttt{coupon\_code} & TEXT & Código de cupón aplicado \\
  \texttt{discount\_amount} & DECIMAL(10,2) & Total de descuento \\
  \texttt{source} & TEXT & direct / booking / airbnb \\
  \texttt{locale} & TEXT & Idioma del proceso (es/en/\ldots) \\
  \texttt{pre\_checkin\_completed} & BOOLEAN & ¿Pre-check-in completado? \\
  \texttt{pre\_checkin\_data} & JSONB & Datos SES.HOSPEDERÍA (cifrar) \\
  \texttt{created\_at} & TIMESTAMPTZ & Timestamp de creación \\
  \bottomrule
\end{longtable}

\subsection{Tabla blocked\_dates}

\begin{longtable}{L{3.5cm} L{3.5cm} L{5.5cm}}
  \toprule
  \rowcolor{AzulejoClaro}
  \textbf{Columna} & \textbf{Tipo} & \textbf{Descripción} \\
  \midrule
  \texttt{date} & DATE UNIQUE & Fecha bloqueada \\
  \texttt{reason} & TEXT & Motivo del bloqueo \\
  \texttt{source} & TEXT & manual / ical\_booking / ical\_airbnb / booking \\
  \texttt{booking\_ref} & TEXT & Referencia de reserva (si aplica) \\
  \bottomrule
\end{longtable}

\subsection{Tabla reviews}

Almacena reseñas con soporte multilingüe completo: columnas
\texttt{text\_es}, \texttt{text\_en}, \texttt{text\_fr}, \texttt{text\_de},
\texttt{text\_it}, \texttt{text\_pt}. Solo se muestran públicamente las que tienen
\texttt{approved = true}.

\subsection{Tabla coupons}

\begin{longtable}{L{3.5cm} L{3.5cm} L{5.5cm}}
  \toprule
  \rowcolor{AzulejoClaro}
  \textbf{Columna} & \textbf{Tipo} & \textbf{Descripción} \\
  \midrule
  \texttt{code} & TEXT UNIQUE & Código en mayúsculas \\
  \texttt{discount\_type} & TEXT & \texttt{percent} o \texttt{fixed} \\
  \texttt{discount\_value} & DECIMAL(10,2) & Valor del descuento \\
  \texttt{valid\_from} & DATE & Inicio de validez \\
  \texttt{valid\_until} & DATE & Fin de validez \\
  \texttt{max\_uses} & INTEGER & Límite de usos (\texttt{null} = ilimitado) \\
  \texttt{uses\_count} & INTEGER & Usos consumidos \\
  \texttt{active} & BOOLEAN & Estado activo/inactivo \\
  \bottomrule
\end{longtable}

\section{Políticas RLS}

\begin{lstlisting}[style=sqlStyle]
-- Solo resenas aprobadas son publicas
CREATE POLICY "Public read approved reviews"
  ON reviews FOR SELECT USING (approved = true);

-- Disponibilidad publica (calendario)
CREATE POLICY "Public read blocked dates"
  ON blocked_dates FOR SELECT USING (true);

-- Insercion publica de reservas (pre-pago)
CREATE POLICY "Public insert bookings"
  ON bookings FOR INSERT WITH CHECK (true);
\end{lstlisting}

%=============================================================
\chapter{Motor de Precios}
%=============================================================

\section{Fórmula matemática}

Sea $n$ el número de noches, $h$ el número de huéspedes, $p$ el porcentaje de
descuento directo y $d_c$ el descuento por cupón:

\begin{align}
  S &= n \times 142                     & \text{(subtotal)} \\
  L &= 60                                & \text{(tarifa limpieza)} \\
  T &= h \times n \times 1{,}50          & \text{(tasa turística)} \\
  D &= \left\lfloor S \times \frac{p}{100} \right\rfloor + d_c & \text{(descuento total)} \\
  \text{Total} &= S + L + T + U - D      & \text{($U$: upsells)} \\
  P_{\text{OTA}} &= \lceil \text{Total} \times 1{,}10 \rceil & \text{(estimación Booking)} \\
  \text{Ahorro} &= P_{\text{OTA}} - \text{Total}               & \text{(visible al usuario)}
\end{align}

\section{Multiplicadores de temporada}

\begin{longtable}{L{6.5cm} C{2.5cm} R{2cm} R{2cm}}
  \toprule
  \rowcolor{AzulejoClaro}
  \textbf{Período} & \textbf{Multiplicador} & \textbf{€/noche} & \textbf{Est. OTA} \\
  \midrule
  Semana Santa (25 mar -- 10 abr) & $\times 1{,}50$ & \euro{}213 & \euro{}234 \\
  Feria de Abril (15 -- 30 abr)   & $\times 1{,}50$ & \euro{}213 & \euro{}234 \\
  Verano (jun -- ago)             & $\times 1{,}25$ & \euro{}178 & \euro{}196 \\
  Navidad / NYE (22 -- 31 dic)    & $\times 1{,}30$ & \euro{}185 & \euro{}204 \\
  Primavera / Otoño               & $\times 1{,}00$ & \euro{}142 & \euro{}156 \\
  Invierno (ene -- feb)           & $\times 0{,}85$ & \euro{}121 & \euro{}133 \\
  \bottomrule
\end{longtable}

\section{Extras disponibles (Upsells)}

\begin{longtable}{L{4.5cm} C{1.8cm} L{6cm}}
  \toprule
  \rowcolor{AzulejoClaro}
  \textbf{Extra} & \textbf{Precio} & \textbf{Descripción} \\
  \midrule
  Traslado aeropuerto SVQ & \euro{}35 & Transfer desde/hasta el aeropuerto \\
  Pack bienvenida local   & \euro{}25 & Productos artesanales sevillanos \\
  Late check-out          & \euro{}20 & Salida hasta las 14:00 h \\
  Cuna de viaje           & \euro{}15 & Para bebés y niños pequeños \\
  Pack romántico          & \euro{}45 & Vino local + flores + decoración \\
  \bottomrule
\end{longtable}

%=============================================================
\chapter{Internacionalización}
%=============================================================

\section{Configuración de locales}

\begin{lstlisting}[style=tsStyle]
export const locales = ['es', 'en', 'fr', 'de', 'it', 'pt'] as const;
export type Locale = typeof locales[number];
export const defaultLocale: Locale = 'es';

export const localeNames: Record<Locale, string> = {
  es: 'Español', en: 'English', fr: 'Français',
  de: 'Deutsch', it: 'Italiano', pt: 'Português',
};
\end{lstlisting}

\section{Sitemap multilingüe}

El sitemap dinámico genera \textbf{30 URLs} (6 locales × 5 páginas) con
\texttt{alternates.languages} completo en cada entrada, habilitando las etiquetas
\texttt{hreflang} para posicionamiento en los motores de búsqueda de cada país.

\section{Localización de fechas}

\texttt{date-fns} provee locales nativos para los 6 idiomas. El \texttt{DayPicker}
y todas las fechas formateadas en la interfaz usan el locale activo:

\begin{lstlisting}[style=tsStyle]
// Resultado de format(new Date('2026-04-20'), 'PPP', { locale })
// es: "20 de abril de 2026"
// en: "April 20th, 2026"
// de: "20. April 2026"
// fr: "20 avril 2026"
\end{lstlisting}

%=============================================================
\chapter{Seguridad}
%=============================================================

\begin{longtable}{L{4.5cm} L{8cm}}
  \toprule
  \rowcolor{AzulejoClaro}
  \textbf{Vector de ataque} & \textbf{Mitigación implementada} \\
  \midrule
  SQL Injection & Supabase usa queries parametrizadas. Nunca interpolación de strings en SQL. \\
  \addlinespace
  XSS & React escapa HTML por defecto. No se usa \texttt{dangerouslySetInnerHTML}. \\
  \addlinespace
  CSRF & Stripe webhooks verificados con HMAC-SHA256. Cookie admin con SameSite=Strict. \\
  \addlinespace
  Price Tampering & El precio se calcula 100\% en el servidor. El cliente solo envía fechas y huéspedes. \\
  \addlinespace
  Acceso admin & Cookie httpOnly — inaccesible desde JavaScript. \\
  \addlinespace
  Secrets exposure & Las claves privadas NUNCA tienen el prefijo \texttt{NEXT\_PUBLIC\_}. \\
  \addlinespace
  Coupon abuse & Validación en servidor + \texttt{uses\_count} atómico en BD. \\
  \addlinespace
  Stripe replay & Webhook verifica timestamp (tolerancia 5 min). \\
  \addlinespace
  Rate limiting & Recomendado Upstash Ratelimit en producción. \\
  \bottomrule
\end{longtable}

\section{Headers de seguridad recomendados}

\begin{lstlisting}[style=tsStyle]
// next.config.mjs
async headers() {
  return [{
    source: '/(.*)',
    headers: [
      { key: 'X-Content-Type-Options', value: 'nosniff' },
      { key: 'X-Frame-Options', value: 'DENY' },
      { key: 'X-XSS-Protection', value: '1; mode=block' },
      { key: 'Referrer-Policy', value: 'strict-origin-when-cross-origin' },
      { key: 'Permissions-Policy',
        value: 'camera=(), microphone=(), geolocation=()' },
    ]
  }];
}
\end{lstlisting}

%=============================================================
\chapter{Despliegue}
%=============================================================

\section{Vercel (recomendado)}

\begin{lstlisting}[style=bashStyle]
npm i -g vercel
cd esencia-sevilla
vercel --prod
\end{lstlisting}

Configuración en el dashboard de Vercel:
\begin{itemize}
  \item Framework: \textbf{Next.js} (autodetectado)
  \item Build command: \texttt{npm run build}
  \item Node.js version: \textbf{20.x}
  \item Root directory: \texttt{esencia-sevilla}
\end{itemize}

\section{Self-hosted con NGINX}

\begin{lstlisting}[style=bashStyle]
git clone https://github.com/danimap27/esencia-sevilla
cd esencia-sevilla && npm install
cp .env.example .env.local  # rellenar con claves reales
npm run build && npm start  # puerto 3000
\end{lstlisting}

Configuración de NGINX como reverse proxy con SSL (Let's Encrypt):

\begin{lstlisting}[style=bashStyle]
server {
  listen 443 ssl http2;
  server_name esenciasevilla.com;
  ssl_certificate /etc/letsencrypt/live/.../fullchain.pem;
  ssl_certificate_key /etc/letsencrypt/live/.../privkey.pem;

  location / {
    proxy_pass http://localhost:3000;
    proxy_set_header Host $host;
    proxy_set_header X-Real-IP $remote_addr;
    proxy_set_header X-Forwarded-Proto $scheme;
  }
}
\end{lstlisting}

\section{Checklist post-despliegue}

\begin{itemize}
  \item[$\square$] Verificar propiedad en Google Search Console; enviar sitemap.xml
  \item[$\square$] Configurar webhook Stripe: \texttt{https://dominio.com/api/webhooks/stripe}
  \item[$\square$] Activar cifrado de datos en Supabase dashboard
  \item[$\square$] Probar reserva completa con tarjeta test de Stripe
  \item[$\square$] Comprobar ChatWidget con API key de OpenRouter
  \item[$\square$] Verificar sincronización iCal con reserva de prueba en Booking.com
  \item[$\square$] Test de Core Web Vitals en PageSpeed Insights
  \item[$\square$] Verificar hreflang en \url{https://www.hreflang.org/checker/}
\end{itemize}

%=============================================================
\chapter{Cumplimiento Legal}
%=============================================================

\begin{longtable}{L{5.5cm} C{1.5cm} L{5.5cm}}
  \toprule
  \rowcolor{AzulejoClaro}
  \textbf{Requisito} & \textbf{Estado} & \textbf{Detalle} \\
  \midrule
  Nº Reg. Turístico visible & \cmark & AT/SE/03584 en header, hero, footer \\
  Política de Privacidad RGPD & \cmark & \texttt{/[locale]/privacidad} \\
  Términos y Condiciones & \pending & Pendiente: \texttt{/[locale]/terminos} \\
  Banner de cookies RGPD & \pending & Pendiente implementación \\
  Pre-check-in digital & \cmark & Portal de huéspedes \\
  Datos bancarios & \cmark & Solo en Stripe PCI-DSS \\
  Derechos ARCO & \cmark & \texttt{hola@esenciasevilla.com} \\
  SES.HOSPEDERÍA & \pending & Pendiente integración API \\
  \bottomrule
\end{longtable}

Bases legales de los tratamientos de datos:

\begin{longtable}{L{5cm} L{4cm} L{3.5cm}}
  \toprule
  \rowcolor{AzulejoClaro}
  \textbf{Tratamiento} & \textbf{Base legal} & \textbf{Artículo RGPD} \\
  \midrule
  Gestión de reservas & Contrato & Art. 6.1.b \\
  Registro de viajeros & Oblig. legal & Art. 6.1.c \\
  Comunicaciones & Contrato & Art. 6.1.b \\
  Marketing & Consentimiento & Art. 6.1.a \\
  Analytics & Interés legítimo & Art. 6.1.f \\
  \bottomrule
\end{longtable}

%=============================================================
\chapter{Trabajo Futuro}
%=============================================================

Este capítulo describe la hoja de ruta completa del producto: mejoras técnicas
pendientes, nuevas funcionalidades, optimizaciones de negocio y visión estratégica
a largo plazo.

%-------------------------------------------------------------
\section{Corto plazo — v1.1 (1–3 meses)}
%-------------------------------------------------------------

\subsection{Páginas pendientes de implementar}

\paragraph{Términos y Condiciones (\texttt{/[locale]/terminos})}
Requerida por RGPD y la normativa turística andaluza. Debe cubrir: política de
cancelación (no reembolsable en reservas directas), responsabilidades del arrendador
y el arrendatario, legislación aplicable (Código Civil español y normativa turística
de la Junta de Andalucía), y jurisdicción competente.

\paragraph{Confirmación de reserva (\texttt{/[locale]/reserva-confirmada})}
Página de éxito post-pago, accesible con \texttt{?ref=ES-XXXX\&session\_id=...}.
Contenido: resumen de la reserva, próximos pasos (pre-check-in, instrucciones de
llegada), código de la caja de llaves (visible solo tras confirmar identidad),
descarga de comprobante PDF, y acceso al portal de huéspedes.

\paragraph{Blog (\texttt{/[locale]/blog})}
Blog turístico sobre Sevilla con objetivo de captación SEO. Los artículos se definen
en \texttt{src/data/blog-posts.ts} con JSON-LD \texttt{Article} para cada entrada.
Contenido planificado:
\begin{itemize}
  \item ``Los 10 mejores espectáculos de flamenco en Sevilla'' — keyword: \emph{flamenco sevilla}
  \item ``Guía completa de la Feria de Abril 2027'' — keyword: \emph{feria de abril}
  \item ``Qué comer en Sevilla: las mejores tapas del centro'' — keyword: \emph{tapas sevilla}
  \item ``Semana Santa en Sevilla: cómo verla sin perderse nada''
\end{itemize}

\paragraph{Portal de huéspedes / Pre-check-in (\texttt{/[locale]/portal/[bookingRef]})}
Formulario de pre-check-in con los datos exigidos por la
\textbf{Ley Orgánica 4/2015 de Seguridad Ciudadana} para el sistema
\textbf{SES.HOSPEDERÍA}:
\begin{itemize}
  \item DNI/NIE/pasaporte (número, tipo, país de expedición, fecha de expedición)
  \item Fecha de nacimiento
  \item Nacionalidad
  \item Dirección de procedencia
\end{itemize}
Los datos se cifran en tránsito, se almacenan en \texttt{bookings.pre\_checkin\_data}
(JSONB cifrado en Supabase Vault), y opcionalmente se envían a SES.HOSPEDERÍA
vía su API REST.

\subsection{Banner de cookies RGPD}

Implementar un banner conforme al RGPD con tres niveles:
\begin{enumerate}
  \item \textbf{Primer nivel}: Aceptar todo / Rechazar todo / Gestionar preferencias.
  \item \textbf{Segundo nivel}: Panel de preferencias por categoría
        (Esenciales, Analytics, Marketing).
  \item \textbf{Persistencia}: Decisión en \texttt{localStorage}; scripts de
        analytics solo se cargan tras aceptación.
\end{enumerate}

\begin{warnbox}{Aviso legal}
Según la Guía sobre el uso de cookies de la AEPD (2023), el simple aviso sin
opción de rechazo no es suficiente. Es obligatorio ofrecer la opción de rechazar
las cookies no esenciales con la misma facilidad con la que se aceptan.
\end{warnbox}

\subsection{Fotos del apartamento}

Añadir las fotos reales a \texttt{public/fotos/}:
\begin{itemize}
  \item Mínimo 12 fotos: dormitorios, salón, cocina, baño, terraza, exterior
  \item Formato preferido: \textbf{WebP} (calidad 85) con fallback JPG
  \item Dimensiones: 1200×800 px (mínimo) para la galería, 800×533 para mobile
  \item OG image: \texttt{public/og-image.jpg} (1200×630 px)
  \item Script de conversión: \texttt{npx sharp-cli --format webp --quality 85 fotos/\*.jpg}
\end{itemize}

\subsection{Emails HTML con diseño de marca}

Las plantillas de confirmación deben incluir:
\begin{itemize}
  \item Logo y colores de Esencia Sevilla (terracota + crema)
  \item Resumen de la reserva con desglose de precios
  \item Mapa de ubicación embebido (imagen estática de Google Maps Static API)
  \item Botón CTA para completar el pre-check-in
  \item Links a la guía turística y el chatbot
  \item Footer RGPD con enlace a la política de privacidad
\end{itemize}

\subsection{Rate limiting en las APIs}

\begin{lstlisting}[style=tsStyle]
// Con Upstash Ratelimit
import { Ratelimit } from '@upstash/ratelimit';
const ratelimit = new Ratelimit({
  redis: Redis.fromEnv(),
  limiter: Ratelimit.slidingWindow(10, '60 s'),
});
const { success } = await ratelimit.limit(ip);
if (!success) return NextResponse.json(
  { error: 'Too many requests' }, { status: 429 });
\end{lstlisting}

Límites recomendados:
\begin{itemize}
  \item \texttt{/api/chat}: 10 peticiones/minuto por IP
  \item \texttt{/api/bookings/checkout}: 5 peticiones/minuto por IP
  \item \texttt{/api/admin/auth}: 3 intentos/minuto (proteger contra bruteforce)
\end{itemize}

%-------------------------------------------------------------
\section{Medio plazo — v1.2 (3–6 meses)}
%-------------------------------------------------------------

\subsection{Precios dinámicos en el calendario}

Mostrar el precio por noche de cada día directamente en el DayPicker, usando
\texttt{getDynamicPrice(day)} que aplica el multiplicador de temporada:

\begin{lstlisting}[style=tsStyle]
const renderDay = (day: Date) => (
  <div className="flex flex-col items-center">
    <span>{day.getDate()}</span>
    <span className="text-[10px] text-terracota-600 font-medium">
      {formatPrice(getDynamicPrice(day))}
    </span>
  </div>
);
\end{lstlisting}

El cálculo del total deberá acumular el precio dinámico de cada noche individualmente
en lugar de multiplicar el precio base por el número de noches.

\subsection{Sistema de reseñas con moderación}

\paragraph{Solicitud automática} 3 días después del check-out, enviar email con
enlace único a \texttt{/[locale]/opinar/[bookingRef]?token=...}

\paragraph{Formulario de reseña} Valoración 1--5 estrellas, texto libre (\textless 500 chars),
idioma autodetectado, token único anti-spam.

\paragraph{Moderación en admin} Tabla con reseñas pendientes, botones Aprobar/Rechazar.

\paragraph{Integración con Google} Tras aprobación, sugerir al huésped que comparta la
reseña en Google Business Profile (enlace directo).

\subsection{Portal de huéspedes completo}

\begin{description}
  \item[Pre-check-in] Datos SES.HOSPEDERÍA como se describe en §11.1.4.
  \item[Información de estancia] Código de la caja de llaves, WiFi, instrucciones de entrada.
  \item[Guía del apartamento] Cómo funciona el climatizador, lavadora, SmartTV.
  \item[Chatbot contextual] El chatbot sabe el nombre del huésped y sus fechas de estancia.
  \item[Checkout digital] Checklist de salida para minimizar las reclamaciones.
\end{description}

\subsection{Suite de testing automatizado}

\paragraph{Unit tests (Vitest)}

\begin{lstlisting}[style=tsStyle]
describe('calculatePrice', () => {
  it('applies 10% direct discount correctly', () => {
    const { discount, total } = calculatePrice(
      new Date('2026-06-15'), new Date('2026-06-20'), 2, 0, 10
    );
    expect(discount).toBe(Math.floor(5 * 142 * 0.1));
  });
  it('estimates Booking.com at +10%', () => {
    const { total, bookingPrice } = calculatePrice(
      new Date('2026-06-15'), new Date('2026-06-20'), 2
    );
    expect(bookingPrice).toBe(Math.round(total * 1.1));
  });
});
\end{lstlisting}

\paragraph{E2E con Playwright}
\begin{itemize}
  \item Flujo completo de reserva con tarjeta de test Stripe (\texttt{4242 4242 4242 4242})
  \item Cambio de idioma y verificación de traducciones
  \item Login admin + creación de bloqueo de fecha
  \item Chat widget: envío de mensaje y recepción de respuesta IA
\end{itemize}

\paragraph{GitHub Actions CI}
\begin{lstlisting}[style=bashStyle]
name: Tests CI
on: [push, pull_request]
jobs:
  test:
    runs-on: ubuntu-latest
    steps:
      - uses: actions/checkout@v4
      - run: npm ci && npm test && npm run test:e2e
\end{lstlisting}

%-------------------------------------------------------------
\section{Largo plazo — v2.0 (6–18 meses)}
%-------------------------------------------------------------

\subsection{Aplicación móvil nativa (React Native / Expo)}

App para huéspedes disponible en App Store y Google Play:

\begin{itemize}
  \item Llave digital NFC o QR para acceso sin caja de llaves
  \item Chat en tiempo real con el propietario (Supabase Realtime)
  \item Guía turística offline con mapas descargados
  \item Notificaciones push: recordatorio check-in, código de acceso, solicitud de reseña
  \item Historial de estancias anteriores
\end{itemize}

\textbf{Stack:} Expo SDK 51 + React Native, Supabase Realtime, Expo Notifications.

\subsection{Yield Management — Revenue Management}

Sistema de precios dinámicos basado en demanda real:

\begin{equation}
  P_{\text{óptimo}}(d) = P_{\text{base}} \cdot M_{\text{temporada}}(d)
    \cdot M_{\text{urgencia}}(\Delta t) \cdot M_{\text{ocupación}}(\rho)
\end{equation}

Donde $\Delta t$ es el número de días hasta la fecha y $\rho$ es la tasa de
ocupación del período.

\begin{longtable}{L{4cm} C{3cm} L{5.5cm}}
  \toprule
  \rowcolor{AzulejoClaro}
  \textbf{Señal} & \textbf{Multiplicador} & \textbf{Lógica} \\
  \midrule
  Lead time < 7 días & $\times 1{,}20$ & Urgencia; el viajero ya está en Sevilla \\
  Lead time < 30 días & $\times 1{,}10$ & Demanda próxima \\
  Ocupación > 80\% & $\times 1{,}15$ & Semana muy demandada \\
  Ocupación < 30\% & $\times 0{,}90$ & Incentivar la reserva \\
  \bottomrule
\end{longtable}

\textbf{Dashboard de Revenue Management} con:
\begin{itemize}
  \item Calendario de precios editable día a día
  \item Gráfico ingresos vs ocupación por mes
  \item RevPAN (Revenue per Available Night), ADR, tasa de ocupación
  \item Comparativa con precios de la competencia en Booking.com
\end{itemize}

\subsection{Soporte multi-propiedad}

Escalar el sistema para gestionar una cartera de apartamentos:

\begin{lstlisting}[style=sqlStyle]
CREATE TABLE properties (
  id UUID PRIMARY KEY DEFAULT gen_random_uuid(),
  slug TEXT UNIQUE NOT NULL,
  owner_id UUID REFERENCES auth.users(id),
  name TEXT NOT NULL,
  address TEXT NOT NULL,
  base_price DECIMAL(10,2) NOT NULL,
  max_guests INTEGER NOT NULL,
  -- ...
);

ALTER TABLE bookings ADD COLUMN property_id UUID
  REFERENCES properties(id) ON DELETE RESTRICT;
ALTER TABLE blocked_dates ADD COLUMN property_id UUID
  REFERENCES properties(id) ON DELETE CASCADE;
\end{lstlisting}

\textbf{Routing multi-propiedad:} \texttt{/[locale]/[property]\#reservar}

\textbf{Autenticación multi-rol:} migrar de cookie simple a Supabase Auth con roles
(propietario, gestor, personal de limpieza).

\subsection{Integraciones con Channel Managers}

\begin{longtable}{L{3cm} L{3cm} L{7cm}}
  \toprule
  \rowcolor{AzulejoClaro}
  \textbf{Sistema} & \textbf{Integración} & \textbf{Beneficio} \\
  \midrule
  Lodgify & API REST & Gestión unificada de todos los canales \\
  Hostaway & API REST & Sincronización bidireccional en tiempo real \\
  Beds24 & iCal + API & Precio dinámico en todos los canales simultáneamente \\
  Guesty & Webhook & Automatización operacional completa \\
  Smoobu & API REST & Solución integrada para propietarios pequeños \\
  \bottomrule
\end{longtable}

\subsection{Casa inteligente (IoT)}

\begin{description}
  \item[Cerradura inteligente] Código temporal generado automáticamente para cada
    huésped, válido desde el check-in hasta el check-out, enviado por email y
    visible en el portal. Stack recomendado: \textbf{Nuki Smart Lock 4.0} con API REST.

  \item[Control climático] Preconfigurar temperatura 1 hora antes del check-in.
    Ahorro energético en períodos sin huésped. Integración: \textbf{Ecobee} o
    \textbf{Google Nest} API.

  \item[Sensor de ruido] Alerta al propietario si se detectan niveles >75 dB en
    horario de silencio (22:00--08:00). Producto: \textbf{Minut} o \textbf{NoiseAware}.

  \item[Detector de CO\textsubscript{2}] Notificación automática si los niveles
    suponen riesgo para la salud. Integración con el sistema de notificaciones
    existente (WhatsApp + email).
\end{description}

\subsection{Chatbot IA avanzado con RAG}

Enriquecer el chatbot con Retrieval-Augmented Generation:

\begin{description}
  \item[Base de conocimiento dinámica] Menús de restaurantes actualizados semanalmente,
    horarios reales de museos y atracciones (scraping), eventos locales en tiempo real
    (API del Ayuntamiento de Sevilla), previsión del tiempo para la semana de estancia.

  \item[Historial persistente] Guardar conversaciones en Supabase para que el bot
    recuerde el contexto entre sesiones del mismo huésped.

  \item[Escalado a humano] Si la confianza de la respuesta es baja, transferir
    automáticamente a WhatsApp del propietario con el historial completo.

  \item[Soporte multimodal] Permitir envío de fotos para soporte visual
    (``¿Cómo funciona este electrodoméstico?'').
\end{description}

\subsection{Dashboard de Business Intelligence}

\begin{description}
  \item[Ingresos] RevPAN, ADR, ingresos mes/trimestre/año, proyección anual.
  \item[Ocupación] Tasa de ocupación, mapa de calor calendario, lead time medio.
  \item[Canales] Reservas directas vs OTAs, comisiones pagadas, ahorro total acumulado.
  \item[Huéspedes] Distribución geográfica (mapa), idioma más frecuente, ratio de repetición, NPS.
  \item[Valoraciones] Evolución de rating en el tiempo, distribución por fuente.
  \item[Exportación] Informes PDF para declaraciones de impuestos (IAE, IVA, IRPF modelo 415).
\end{description}

\subsection{Automatización de operaciones}

\begin{longtable}{L{6cm} L{6.5cm}}
  \toprule
  \rowcolor{AzulejoClaro}
  \textbf{Trigger} & \textbf{Acción automática} \\
  \midrule
  7 días antes check-in & Email: instrucciones de llegada + link pre-check-in \\
  1 día antes check-in & SMS + email: código de la caja de llaves \\
  Día check-in (16:00) & WhatsApp bienvenida + link guía turística \\
  Día check-out (09:00) & Recordatorio horario de salida \\
  3 días post check-out & Solicitud de reseña por email \\
  10 días sin reseña & Segundo recordatorio de reseña \\
  Fin de mes & Resumen de ingresos por email al propietario \\
  \bottomrule
\end{longtable}

\textbf{Stack:} Supabase Edge Functions + pg\_cron, o servicio externo (Zapier / Make.com).

\subsection{Marketplace de experiencias locales}

Ofrecer actividades curadas como upsells adicionales:

\begin{longtable}{L{6.5cm} C{2cm} C{2cm}}
  \toprule
  \rowcolor{AzulejoClaro}
  \textbf{Experiencia} & \textbf{Precio} & \textbf{Comisión} \\
  \midrule
  Espectáculo flamenco (2 pax) & \euro{}85 & 20\% \\
  Tour de tapas y vinos (4 pax) & \euro{}120 & 15\% \\
  Ruta en bici por el Guadalquivir & \euro{}40/pax & 18\% \\
  Paddle surf en el río & \euro{}35/pax & 20\% \\
  Clase de cocina sevillana & \euro{}65/pax & 15\% \\
  Visita privada al Alcázar & \euro{}95 & 20\% \\
  \bottomrule
\end{longtable}

\textbf{Integración:} GetYourGuide Affiliate API o Airbnb Experiences para inventario en tiempo real.

\subsection{App para el personal de limpieza}

Aplicación simplificada con:
\begin{itemize}
  \item Calendario de próximas limpiezas (check-outs pendientes)
  \item Checklist de limpieza con confirmación fotográfica por estancia
  \item Gestión de inventario (productos de limpieza, ropa de cama, toallas)
  \item Reporte de incidencias con foto, descripción y nivel de urgencia
\end{itemize}

%-------------------------------------------------------------
\section{Mejoras técnicas transversales}
%-------------------------------------------------------------

\subsection{Migración a Supabase Auth}

Sustituir la cookie simple del admin por \textbf{Supabase Auth} con:
\begin{itemize}
  \item JWT tokens con refresh automático
  \item Múltiples usuarios con roles (admin, gestor, limpieza)
  \item 2FA/MFA para el propietario (TOTP o SMS)
  \item Magic Links para acceso sin contraseña
\end{itemize}

\subsection{Internacionalización extendida}

\begin{longtable}{L{4cm} C{2.5cm} L{6cm}}
  \toprule
  \rowcolor{AzulejoClaro}
  \textbf{Idioma} & \textbf{Prioridad} & \textbf{Justificación} \\
  \midrule
  Japonés (ja) & Alta & Turismo japonés creciente en Andalucía \\
  Chino simplificado (zh) & Alta & Mayor mercado turístico del mundo \\
  Árabe (ar) + RTL & Media & Turismo del Golfo Pérsico en crecimiento \\
  Ruso (ru) & Media & Frecuente en turismo de Europa del Este \\
  Holandés (nl) & Media & Alta incidencia en turismo español \\
  \bottomrule
\end{longtable}

También: detección de idioma por \textbf{GeoIP} (no solo Accept-Language),
y \textbf{conversión de moneda} para mostrar precios en GBP, USD, CHF manteniendo
el cobro en EUR.

\subsection{Accesibilidad (WCAG 2.2 AA)}

\begin{itemize}
  \item Navegación completa por teclado (Tab, Enter, Space, Escape, flechas)
  \item Skip-to-content links
  \item ARIA labels en todos los elementos interactivos sin texto visible
  \item Contraste de color $\geq$ 4,5:1 en todos los textos
  \item Alt text descriptivo en todas las imágenes
  \item Respeto a \texttt{prefers-reduced-motion} para animaciones
  \item Verificación con Axe DevTools + Lighthouse Accessibility audit
\end{itemize}

\subsection{Observabilidad y monitorización}

\begin{description}
  \item[Error tracking] \textbf{Sentry} para capturar errores en producción con
    contexto completo (stack trace, usuario, request).

  \item[Uptime monitoring] \textbf{BetterUptime} o \textbf{UptimeRobot} con
    alertas por email/SMS si la web cae.

  \item[Logs estructurados] \textbf{Pino} logger en API routes, exportado a
    Logtail o Datadog.

  \item[Health check] \texttt{GET /api/health} que verifica la conexión a
    Supabase, Stripe y devuelve la versión y uptime.
\end{description}

\subsection{Edge Runtime para disponibilidad}

\begin{lstlisting}[style=tsStyle]
// app/api/availability/route.ts
export const runtime = 'edge';  // latencia < 50ms desde cualquier region
\end{lstlisting}

\subsection{CDN para imágenes}

Migrar las fotos del apartamento a \textbf{Cloudinary} o \textbf{ImageKit}:
\begin{itemize}
  \item Conversión automática a AVIF/WebP según soporte del navegador
  \item Resize bajo demanda sin reprocessing en el servidor
  \item CDN global con edge nodes en Europa para baja latencia
  \item LQIP (Low Quality Image Placeholders) generados automáticamente
\end{itemize}

%=============================================================
% Apéndices
%=============================================================

\appendix

\chapter{Variables de entorno completas}

\begin{lstlisting}[style=bashStyle]
# URL del sitio
NEXT_PUBLIC_SITE_URL=https://esenciasevilla.com

# Supabase
NEXT_PUBLIC_SUPABASE_URL=https://xxx.supabase.co
NEXT_PUBLIC_SUPABASE_ANON_KEY=eyJ...
SUPABASE_SERVICE_ROLE_KEY=eyJ...

# Stripe
STRIPE_SECRET_KEY=sk_live_...
NEXT_PUBLIC_STRIPE_PUBLISHABLE_KEY=pk_live_...
STRIPE_WEBHOOK_SECRET=whsec_...

# IA
OPENROUTER_API_KEY=sk-or-v1-...

# Comunicaciones
RESEND_API_KEY=re_...
TWILIO_ACCOUNT_SID=AC...
TWILIO_AUTH_TOKEN=...

# iCal sync OTAs
ICAL_BOOKING_URL=https://admin.booking.com/hotel/.../ical.html?...
ICAL_AIRBNB_URL=https://www.airbnb.com/calendar/ical/...

# Analytics
NEXT_PUBLIC_GA_MEASUREMENT_ID=G-XXXXXXXXXX
NEXT_PUBLIC_META_PIXEL_ID=...

# Panel admin
ADMIN_PASSWORD=password_muy_segura_min12_chars
\end{lstlisting}

\chapter{Schema SQL completo de Supabase}

\begin{lstlisting}[style=sqlStyle]
-- Tabla de reservas
CREATE TABLE bookings (
  id UUID PRIMARY KEY DEFAULT gen_random_uuid(),
  booking_ref TEXT UNIQUE NOT NULL,
  check_in DATE NOT NULL, check_out DATE NOT NULL,
  guests INTEGER NOT NULL DEFAULT 1,
  guest_name TEXT NOT NULL, guest_email TEXT NOT NULL,
  guest_phone TEXT,
  total_amount DECIMAL(10,2) NOT NULL,
  status TEXT NOT NULL DEFAULT 'pending'
    CHECK (status IN ('pending','confirmed','cancelled','completed')),
  payment_intent_id TEXT,
  upsells JSONB DEFAULT '[]',
  coupon_code TEXT, discount_amount DECIMAL(10,2) DEFAULT 0,
  source TEXT DEFAULT 'direct', locale TEXT DEFAULT 'es',
  created_at TIMESTAMPTZ DEFAULT NOW(),
  updated_at TIMESTAMPTZ DEFAULT NOW(),
  pre_checkin_completed BOOLEAN DEFAULT FALSE,
  pre_checkin_data JSONB
);

-- Tabla de fechas bloqueadas
CREATE TABLE blocked_dates (
  id UUID PRIMARY KEY DEFAULT gen_random_uuid(),
  date DATE NOT NULL, reason TEXT,
  source TEXT DEFAULT 'manual', booking_ref TEXT,
  created_at TIMESTAMPTZ DEFAULT NOW()
);
CREATE UNIQUE INDEX blocked_dates_date_idx ON blocked_dates(date);

-- Tabla de resenas
CREATE TABLE reviews (
  id UUID PRIMARY KEY DEFAULT gen_random_uuid(),
  author TEXT NOT NULL, country TEXT, country_flag TEXT,
  rating INTEGER NOT NULL CHECK (rating BETWEEN 1 AND 5),
  review_date DATE NOT NULL, source TEXT DEFAULT 'direct',
  text_es TEXT, text_en TEXT, text_fr TEXT,
  text_de TEXT, text_it TEXT, text_pt TEXT,
  approved BOOLEAN DEFAULT FALSE,
  created_at TIMESTAMPTZ DEFAULT NOW()
);

-- Tabla de cupones
CREATE TABLE coupons (
  id UUID PRIMARY KEY DEFAULT gen_random_uuid(),
  code TEXT UNIQUE NOT NULL,
  discount_type TEXT NOT NULL CHECK (discount_type IN ('percent','fixed')),
  discount_value DECIMAL(10,2) NOT NULL,
  valid_from DATE, valid_until DATE,
  max_uses INTEGER, uses_count INTEGER DEFAULT 0,
  active BOOLEAN DEFAULT TRUE,
  created_at TIMESTAMPTZ DEFAULT NOW()
);

-- RLS
ALTER TABLE bookings ENABLE ROW LEVEL SECURITY;
ALTER TABLE blocked_dates ENABLE ROW LEVEL SECURITY;
ALTER TABLE reviews ENABLE ROW LEVEL SECURITY;
ALTER TABLE coupons ENABLE ROW LEVEL SECURITY;

CREATE POLICY "Public read approved reviews"
  ON reviews FOR SELECT USING (approved = true);
CREATE POLICY "Public read blocked dates"
  ON blocked_dates FOR SELECT USING (true);
CREATE POLICY "Public insert bookings"
  ON bookings FOR INSERT WITH CHECK (true);
\end{lstlisting}

%--- Índices de tablas y figuras --------------------------
\clearpage
\listoftables
\clearpage
\listoffigures

%=============================================================
\end{document}
%=============================================================
